% ============================================================
%  CHUONG 3. HIEN THUC HOA NGHIEN CUU
% ============================================================
\chapter{HIỆN THỰC HÓA NGHIÊN CỨU}
\label{ch:hien-thuc}

Chương này trình bày quá trình hiện thực hóa hệ thống, đi từ mô tả bài toán và
đặc tả yêu cầu tới thiết kế cơ sở dữ liệu, lược đồ use case, kiến trúc hệ thống,
thiết kế lớp, thiết kế luồng xử lý và cuối cùng là chi tiết cài đặt ba thuật toán
lõi.

\section{Mô tả bài toán}
\label{sec:mo-ta-bai-toan}

\subsection{Bối cảnh và phạm vi của bài toán}

Bài toán của đồ án xoay quanh chuỗi công việc hằng ngày của người nội trợ: quyết
định nấu món gì từ nguyên liệu đang có, sắp xếp các bữa ăn cho cả tuần, và tổng
hợp danh sách cần mua. Ba việc này nối tiếp nhau về mặt dữ liệu: đầu ra của việc
trước là đầu vào của việc sau. Đây chính là lý do đồ án chọn hướng tích hợp cả ba
trong một hệ thống thay vì làm ba công cụ rời rạc.

Đầu vào của hệ thống gồm hai nguồn: kho dữ liệu món ăn do quản trị viên duy trì
(mỗi món có công thức, danh sách nguyên liệu kèm định lượng và đơn vị, thông tin
danh mục, vùng miền, độ khó, thời gian và năng lượng theo khẩu phần), và dữ liệu
người dùng cung cấp tại thời điểm sử dụng (tập nguyên liệu đang có, tuần cần lập
thực đơn, lựa chọn vùng miền nếu có). Đầu ra tương ứng gồm ba sản phẩm: danh
sách món đã xếp hạng kèm nhãn tình trạng nguyên liệu, lịch thực đơn 21 suất ăn,
và danh sách đi chợ đã gộp nhóm và cộng dồn khối lượng.

\subsection{Mô hình tổng quát}

Hình~\ref{fig:mo-ta-bai-toan} mô tả bài toán ở mức tổng quát: người nội trợ tương
tác với hệ thống qua ba chức năng nối tiếp nhau, cả ba cùng khai thác một kho dữ
liệu món ăn đã chuẩn hóa về định lượng nguyên liệu.

\begin{figure}[H]
\centering
\resizebox{\textwidth}{!}{%
\begin{tikzpicture}[node distance=1.2cm and 1.2cm]

  \stickman{ND}{0}{0.6}{Người nội trợ}

  \node[procs] (GY) at (4.2,1.2) {(1) Gợi ý món ăn\\theo nguyên liệu sẵn có};
  \node[procs, right=1.0cm of GY] (TD) {(2) Sinh thực đơn tuần\\7 ngày $\times$ 3 bữa,\\chỉnh tay được};
  \node[procs, right=1.0cm of TD] (DC) {(3) Tổng hợp danh sách\\đi chợ, gộp và cộng dồn};

  \node[devgroup, fit=(GY)(TD)(DC), label={[font=\small]above:Hệ thống CookingAdvisor}] (HT) {};

  \node[store, below=1.6cm of TD] (KHO) {Kho dữ liệu:\\món ăn, nguyên liệu,\\định lượng};

  \draw[fl] (GY) -- (TD);
  \draw[fl] (TD) -- (DC);

  \draw[flb] (ND) -- (HT.west);
  \draw[flb] (HT.south) -- (KHO.north);

\end{tikzpicture}%
}
\caption{Sơ đồ mô tả bài toán ở mức tổng quát}
\label{fig:mo-ta-bai-toan}
\end{figure}


Luồng nghiệp vụ chính diễn ra như sau. Người dùng khai báo các nguyên liệu đang
có; hệ thống đối chiếu với kho công thức và trả về danh sách món được xếp hạng
theo mức độ phù hợp, ghi rõ món nào nấu được ngay và món nào còn thiếu nguyên
liệu gì. Khi cần lên kế hoạch cho cả tuần, hệ thống sinh tự động lịch 7 ngày với
3 bữa mỗi ngày, người dùng chỉnh tay từng ô nếu muốn. Cuối cùng, từ thực đơn đã
chốt, hệ thống gộp toàn bộ nguyên liệu của 21 suất ăn thành danh sách đi chợ,
cộng dồn khối lượng theo từng nguyên liệu và đơn vị.

Ba chức năng này tương ứng với ba bài toán con đã phân tích ở
mục~\ref{sec:boi-canh}: truy vấn ngược từ nguyên liệu, xếp lịch có ràng buộc, và
gộp nhóm cùng tổng hợp. Phần còn lại của chương trình bày yêu cầu, thiết kế dữ
liệu và cách cài đặt từng bài toán con đó.

\section{Đặc tả yêu cầu}
\label{sec:dac-ta}

\subsection{Xác định tác nhân}
\label{sec:tac-nhan}

Hệ thống có ba tác nhân, phân biệt theo mức quyền truy cập tăng dần.

\begin{table}[H]
\centering
\caption{Các tác nhân của hệ thống}
\label{tab:tac-nhan}
\begin{tabularx}{\textwidth}{|P{3.0cm}|Y|P{4.8cm}|}
\hline
\bfseries Tác nhân & \bfseries Mô tả & \bfseries Cơ chế nhận diện \\ \hline
\textbf{Khách} & Người truy cập chưa đăng nhập & Không có phiên đăng nhập \\ \hline
\textbf{Người dùng} & Tài khoản đã đăng nhập, vai trò \code{User} &
  Thuộc tính \code{[Authorize]} \\ \hline
\textbf{Quản trị viên} & Tài khoản thuộc vai trò \code{Admin} &
  \code{[Authorize(Roles = "Admin")]} \\ \hline
\end{tabularx}
\end{table}

Quan hệ giữa ba tác nhân là quan hệ kế thừa quyền: Người dùng có toàn bộ quyền
của Khách, Quản trị viên có toàn bộ quyền của Người dùng.

\subsection{Yêu cầu chức năng}
\label{sec:yeu-cau-chuc-nang}

Hệ thống cung cấp 20 chức năng, chia thành ba nhóm theo mức quyền truy cập,
được liệt kê chi tiết ở Bảng~\ref{tab:chuc-nang}.

\begin{longtable}{|P{0.9cm}|P{4.6cm}|P{7.9cm}|}
\caption{Yêu cầu chức năng của hệ thống theo nhóm quyền truy cập}
\label{tab:chuc-nang}\\
\hline
\bfseries Mã & \bfseries Tên chức năng & \bfseries Mô tả \\ \hline
\endfirsthead
\hline
\bfseries Mã & \bfseries Tên chức năng & \bfseries Mô tả \\ \hline
\endhead
\endfoot
\multicolumn{3}{|P{14.24cm}|}{\bfseries Nhóm A. Dành cho Khách (không cần đăng nhập)} \\ \hline
A1 & Xem trang chủ & Hiển thị danh mục, món nổi bật, số liệu tổng quan \\ \hline
A2 & Duyệt danh sách món ăn & Hiển thị dạng lưới, phân trang 9 món mỗi
  trang \\ \hline
A3 & Tìm kiếm theo tên món & Tìm theo chuỗi con trong tên \\ \hline
A4 & Lọc nâng cao & Theo danh mục, vùng miền, độ khó, thời gian nấu tối
  đa \\ \hline
A5 & Sắp xếp kết quả & Theo tên, thời gian nấu, năng lượng tăng hoặc giảm \\ \hline
A6 & Xem chi tiết món ăn & Thông số, nguyên liệu định lượng, các bước
  nấu \\ \hline
A7 & Gợi ý theo nguyên liệu & Chọn nguyên liệu đang có, nhận danh sách món xếp
  hạng \\ \hline
A8 & Đăng ký tài khoản & Tạo tài khoản mới, tự động gán vai trò
  \code{User} \\ \hline
A9 & Đăng nhập, đăng xuất & Xác thực bằng email và mật khẩu \\ \hline
\multicolumn{3}{|P{14.24cm}|}{\bfseries Nhóm B. Dành cho Người dùng đã đăng nhập} \\ \hline
B1 & Sinh thực đơn tuần & Tự động lấp 21 suất ăn, có tùy chọn vùng miền \\ \hline
B2 & Xem lịch thực đơn & Bảng 7 ngày $\times$ 3 bữa kèm tổng năng lượng
  theo ngày \\ \hline
B3 & Sửa thủ công từng bữa & Đổi món hoặc xóa món khỏi một ô \\ \hline
B4 & Xem danh sách thực đơn đã lưu & Liệt kê các thực đơn của riêng người
  dùng \\ \hline
B5 & Tạo danh sách đi chợ & Sinh từ một thực đơn, gộp và cộng dồn nguyên liệu \\ \hline
B6 & Đánh dấu đã mua & Tick từng nguyên liệu, hiển thị tiến độ \\ \hline
B7 & Quản lý món yêu thích & Thêm, bỏ và xem danh sách \\ \hline
\multicolumn{3}{|P{14.24cm}|}{\bfseries Nhóm C. Dành cho Quản trị viên} \\ \hline
C1 & Xem bảng điều khiển & Thống kê số món, nguyên liệu, danh mục \\ \hline
C2 & Quản lý món ăn & Thêm, sửa, xóa; gán nguyên liệu kèm định lượng \\ \hline
C3 & Quản lý nguyên liệu & Thêm, sửa, xóa; có kiểm tra ràng buộc tham chiếu \\ \hline
C4 & Quản lý danh mục & Thêm, sửa, xóa; có kiểm tra ràng buộc tham
  chiếu \\ \hline
\end{longtable}

\subsection{Yêu cầu phi chức năng}
\label{sec:yeu-cau-phi-chuc-nang}

\begin{table}[H]
\centering
\caption{Yêu cầu phi chức năng}
\label{tab:phi-chuc-nang}
\begin{tabularx}{\textwidth}{|P{0.9cm}|P{4.0cm}|Y|}
\hline
\bfseries Mã & \bfseries Yêu cầu & \bfseries Tiêu chí nghiệm thu \\ \hline
N1 & Bảo mật mật khẩu & Mật khẩu lưu dưới dạng băm, tối thiểu 8 ký tự \\ \hline
N2 & Phân quyền & Người dùng thường không truy cập được khu quản trị \\ \hline
N3 & Cách ly dữ liệu & Người dùng chỉ thao tác được trên thực đơn của mình \\ \hline
N4 & Tương thích màn hình & Hiển thị đúng từ 375 px tới 1440 px, không
  cuộn ngang \\ \hline
N5 & Khả năng tiếp cận & Tương phản màu đạt WCAG 2.1 mức AA \\ \hline
N6 & Tính đúng đắn thuật toán & Có kiểm thử đơn vị tự động cho ba thuật
  toán lõi \\ \hline
N7 & Khả năng tái lập & Dựng lại cơ sở dữ liệu bằng một lệnh trên máy khác \\ \hline
\end{tabularx}
\end{table}

\section{Mô hình cơ sở dữ liệu}
\label{sec:mo-hinh-csdl}

\subsection{Sơ đồ quan hệ thực thể}

Cơ sở dữ liệu gồm 9 bảng nghiệp vụ do đồ án tự thiết kế, cộng thêm các bảng do
ASP.NET Core Identity sinh ra để lưu tài khoản và vai trò. Trong sơ đồ ở
Hình~\ref{fig:erd}, thực thể \code{ApplicationUser} đại diện cho bảng người dùng
của Identity đã được mở rộng thêm hai cột riêng của đồ án.

\begin{figure}[H]
\centering
\resizebox{\textwidth}{!}{%
\begin{tikzpicture}[node distance=1.2cm and 1.2cm]

  % --- Cot trai: nguoi dung ---------------------------------
  \node[erdent] (USER) at (0,17) {\textbf{ApplicationUser}
    \nodepart{two} Id PK\\UserName\\Email\\PasswordHash\\FullName\\FamilyName};
  \node[erdent] (PLAN) at (0,11) {\textbf{MenuPlan}
    \nodepart{two} Id PK\\UserId FK\\Name\\WeekStartDate\\CreatedAt};
  \node[erdent] (SLIST) at (0,5) {\textbf{ShoppingList}
    \nodepart{two} Id PK\\MenuPlanId FK UK\\UserId FK\\CreatedAt};

  % --- Cot giua: thuc the giao ------------------------------
  \node[erdent] (FAV) at (6.5,14.5) {\textbf{Favorite}
    \nodepart{two} UserId PK FK\\RecipeId PK FK};
  \node[erdent] (PITEM) at (6.5,9) {\textbf{MenuPlanItem}
    \nodepart{two} Id PK\\MenuPlanId FK\\DayOfWeek\\MealType\\RecipeId FK};
  \node[erdent] (SITEM) at (6.5,3) {\textbf{ShoppingListItem}
    \nodepart{two} Id PK\\ShoppingListId FK\\IngredientId FK\\Quantity\\Unit\\IsPurchased};

  % --- Cot Recipe --------------------------------------------
  \node[erdent] (CAT) at (13,17) {\textbf{Category}
    \nodepart{two} Id PK\\Name};
  \node[erdent] (RECIPE) at (13,10) {\textbf{Recipe}
    \nodepart{two} Id PK\\Name\\Description\\Instructions\\Servings\\PrepMinutes\\CookMinutes\\Difficulty\\Region\\CaloriesPerServing\\ImageUrl\\SuitableMealTypes\\CategoryId FK};

  % --- Cot phai: nguyen lieu -----------------------------------
  \node[erdent] (RI) at (19.5,13.5) {\textbf{RecipeIngredient}
    \nodepart{two} RecipeId PK FK\\IngredientId PK FK\\Quantity\\Unit};
  \node[erdent] (ING) at (19.5,7.5) {\textbf{Ingredient}
    \nodepart{two} Id PK\\Name\\Unit\\Group};

  % --- Quan he: nguoi dung -------------------------------------
  \draw[arr] (USER) -- (PLAN) node[lbl, pos=0.5] {lập};
  \draw[arr] (USER.west) -- ++(-1.8,0) |- (SLIST.west) node[lbl, pos=0.62] {sở hữu};
  \draw[arr] (USER.east) -| (FAV.north) node[lbl, pos=0.75] {đánh dấu};

  % --- Quan he: menu / cho -------------------------------------
  \draw[arr] (PLAN) -- (SLIST) node[lbl, pos=0.5] {sinh ra (1-1)};
  \draw[arr] (PLAN.east) -| (PITEM.north) node[lbl, pos=0.7] {chứa};
  \draw[arr] (SLIST.east) -| (SITEM.north);

  % --- Quan he: recipe / category / ingredient -----------------
  \draw[arr] (CAT) -- (RECIPE) node[lbl, pos=0.5] {phân loại};
  \draw[arr] (RECIPE.west) -| (PITEM.east);
  \draw[arr] (RECIPE.north) -| (FAV.east);
  \draw[arr] (RECIPE.east) -| (RI.west) node[lbl, pos=0.7] {gồm};
  \draw[arr] (ING.north) -- (RI.south) node[lbl, pos=0.5] {được dùng trong};
  \draw[arr] (ING.south) |- ++(0,-4) -| (SITEM.east);

\end{tikzpicture}%
}
\caption{Sơ đồ quan hệ thực thể của cơ sở dữ liệu}
\label{fig:erd}
\end{figure}


\subsection{Mô tả các bảng}
\label{sec:mo-ta-bang}

Bảng~\ref{tab:tom-tat-bang} nêu vai trò và khóa của cả 9 bảng nghiệp vụ. Cấu
trúc chi tiết từng cột kèm kiểu dữ liệu trên SQL Server nằm ở Phụ lục G.

\begin{longtable}{|P{4.5cm}|P{5.2cm}|P{3.7cm}|}
\caption{Vai trò và khóa của các bảng nghiệp vụ}
\label{tab:tom-tat-bang}\\
\hline
\bfseries Bảng & \bfseries Vai trò & \bfseries Khóa \\ \hline
\endfirsthead
\hline
\bfseries Bảng & \bfseries Vai trò & \bfseries Khóa \\ \hline
\endhead
\endfoot
\code{Categories} & Danh mục món ăn, làm bộ lọc trên trang danh sách &
  Khóa chính \code{Id} \\ \hline
\code{Ingredients} & Kho nguyên liệu dùng chung, có cột nhóm để gom dòng khi
  hiển thị danh sách đi chợ & Khóa chính \code{Id} \\ \hline
\code{Recipes} & Bảng trung tâm, lưu công thức cùng thời gian, độ khó, vùng
  miền, năng lượng và các bữa phù hợp & Khóa chính \code{Id}; khóa ngoại
  \code{CategoryId} \\ \hline
\codelong{RecipeIngredients} & Bảng trung gian món ăn - nguyên liệu, mang thêm định
  lượng và đơn vị; là mấu chốt của cả ba thuật toán & Khóa chính ghép
  \codelong{(RecipeId, IngredientId)} \\ \hline
\code{MenuPlans} & Thông tin chung của một thực đơn tuần & Khóa chính \code{Id};
  khóa ngoại \code{UserId} \\ \hline
\code{MenuPlanItems} & Từng ô của lưới thực đơn, mỗi thực đơn đầy đủ gồm 21 bản
  ghi & Khóa chính \code{Id}; khóa ngoại \code{MenuPlanId}, \code{RecipeId} \\ \hline
\code{Favorites} & Dấu yêu thích của người dùng với một món & Khóa chính ghép
  \code{(UserId, RecipeId)} \\ \hline
\code{ShoppingLists} & Thông tin chung của một danh sách đi chợ &
  Khóa chính \code{Id}; \code{MenuPlanId} có ràng buộc duy nhất \\ \hline
\codelong{ShoppingListItems} & Từng dòng nguyên liệu đã gộp và cộng dồn khối lượng &
  Khóa chính \code{Id}; khóa ngoại \codelong{ShoppingListId},
  \code{IngredientId} \\ \hline
\end{longtable}

Ngoài 9 bảng trên, hệ thống kế thừa bộ bảng tài khoản do ASP.NET Core Identity
sinh ra; lớp \code{ApplicationUser} kế thừa lớp người dùng chuẩn và bổ sung hai
cột \code{FullName} (họ tên đầy đủ, bắt buộc) và \code{FamilyName} (tên gia
đình, không bắt buộc), giữ nguyên toàn bộ cơ chế bảo mật sẵn có mà vẫn mở rộng
được dữ liệu riêng của bài toán.

Bốn điểm thiết kế trong lược đồ trên cần được nhấn mạnh. \textbf{Thứ nhất},
trường \code{SuitableMealTypes} dùng kiểu liệt kê dạng cờ với Breakfast = 1,
Lunch = 2, Dinner = 4 nên lưu được tổ hợp nhiều bữa trong một số nguyên (giá trị
6 nghĩa là hợp cả trưa lẫn tối) và phép kiểm tra hợp bữa quy về phép \code{AND}
bit. Kiểu này khác kiểu \code{MealType} của \code{MenuPlanItem}, nơi Breakfast
bằng 0 vì là giá trị đơn; gộp chung một kiểu thì mọi tổ hợp chứa bữa sáng sẽ bị
mất. \textbf{Thứ hai}, bảng \code{RecipeIngredients} dùng khóa chính ghép nên
một nguyên liệu chỉ xuất hiện một lần trong mỗi món, đồng thời mang thêm
\code{Quantity} và \code{Unit}; nếu chỉ lưu quan hệ nhiều - nhiều thuần túy thì
danh sách đi chợ không cộng dồn được khối lượng (mục~\ref{sec:gop-nhom}), và cột
\code{Unit} tách khỏi đơn vị mặc định của bảng \code{Ingredients} cho phép cùng
một nguyên liệu định lượng khác nhau ở hai món. \textbf{Thứ ba}, khóa chính ghép
\code{(UserId, RecipeId)} của bảng \code{Favorites} chặn việc đánh dấu trùng một
món. \textbf{Thứ tư}, ràng buộc duy nhất trên \code{MenuPlanId} của bảng
\code{ShoppingLists} thể hiện quan hệ một - một: mỗi thực đơn chỉ có tối đa một
danh sách đi chợ, nên thao tác tạo lại được cài đặt theo hướng ghi đè, tránh
việc bấm hai lần sinh ra hai danh sách trùng nhau.

\subsection{Thiết kế hành vi xóa}
\label{sec:hanh-vi-xoa}

Hành vi xóa được cấu hình có chủ đích cho từng quan hệ, chia làm hai nhóm.

\begin{table}[H]
\centering
\caption{Hành vi xóa của các quan hệ trong cơ sở dữ liệu}
\label{tab:hanh-vi-xoa}
\begin{tabularx}{\textwidth}{|P{5.0cm}|P{1.9cm}|Y|}
\hline
\bfseries Quan hệ & \bfseries Hành vi & \bfseries Lý do \\ \hline
\code{Recipe} tới \code{RecipeIngredient} & Cascade & Xóa món thì các dòng
  nguyên liệu của nó không còn ý nghĩa \\ \hline
\code{MenuPlan} tới \code{MenuPlanItem} & Cascade & Tương tự, các ô thuộc
  về thực đơn \\ \hline
\code{MenuPlan} tới \code{ShoppingList} & Cascade & Danh sách đi chợ phụ thuộc
  hoàn toàn vào thực đơn \\ \hline
\code{ShoppingList} tới \code{ShoppingListItem} & Cascade & Các dòng thuộc
  về danh sách \\ \hline
\code{ApplicationUser} tới \code{MenuPlan}, \code{Favorite} & Cascade & Xóa tài
  khoản thì dữ liệu cá nhân đi theo \\ \hline
\code{Recipe} tới \code{Favorite} & Cascade & Món bị xóa thì dấu yêu thích
  trên món đó không còn ý nghĩa \\ \hline
\code{Category} tới \code{Recipe} & \textbf{Restrict} & Không cho xóa danh mục
  đang có món sử dụng \\ \hline
\code{Ingredient} tới \code{RecipeIngredient} & \textbf{Restrict} & Không
  cho xóa nguyên liệu đang được dùng \\ \hline
\code{Ingredient} tới \code{ShoppingListItem} & \textbf{Restrict} & Không cho xóa
  nguyên liệu đang nằm trong danh sách đi chợ \\ \hline
\code{ApplicationUser} tới \code{ShoppingList} & \textbf{Restrict} & Danh
  sách đi chợ đã xóa lan truyền theo thực đơn, không cần thêm một đường Cascade
  thứ hai từ tài khoản \\ \hline
\code{Recipe} tới \code{MenuPlanItem} & \textbf{Restrict} & Không cho xóa món
  đang nằm trong thực đơn của người dùng \\ \hline
\end{tabularx}
\end{table}

Nguyên tắc phân biệt là: dữ liệu \textbf{phụ thuộc sự tồn tại} thì xóa lan
truyền, còn dữ liệu \textbf{được tham chiếu} thì chặn xóa để tránh phá vỡ dữ liệu
của người dùng khác. Nếu để \code{Recipe} tới \code{MenuPlanItem} là Cascade,
quản trị viên xóa một món sẽ âm thầm phá vỡ tính toàn vẹn của thực đơn đã lưu của
mọi người dùng.

Hai quan hệ cùng trỏ tới \code{Recipe} được cấu hình khác nhau một cách có chủ
đích: \code{MenuPlanItem} là Restrict còn \code{Favorite} là Cascade. Lý do nằm ở
bản chất dữ liệu. Một suất ăn trong thực đơn là \textbf{sản phẩm có cấu trúc} mà
người dùng đã bỏ công lập và chỉnh sửa; mất một suất là thực đơn thủng một ô.
Ngược lại, dấu yêu thích chỉ là \textbf{liên kết đánh dấu} tới món; khi món không
còn tồn tại thì dấu đó không còn đối tượng để trỏ tới, giữ lại cũng không hiển
thị được gì. Nếu cấu hình Favorite là Restrict, quản trị viên sẽ không thể xóa
một món chỉ vì có người từng đánh dấu yêu thích nó, một ràng buộc gây phiền mà
không bảo vệ được giá trị nào.

\subsection{Khởi tạo dữ liệu mẫu}

Lớp \code{DbInitializer} thực hiện ba việc khi ứng dụng khởi động: áp dụng các
migration còn thiếu, tạo hai vai trò \code{Admin} và \code{User} cùng tài khoản
quản trị mặc định, và nạp dữ liệu mẫu nếu bảng \code{Recipes} còn rỗng. Dữ liệu
mẫu gồm 26 món ăn Việt Nam, 40 nguyên liệu và 10 danh mục.

Điều kiện ``chỉ nạp khi bảng còn rỗng'' bảo đảm thao tác khởi động là \textbf{lũy
đẳng}: chạy lại ứng dụng nhiều lần không tạo ra dữ liệu trùng lặp.

\section{Lược đồ use case}
\label{sec:use-case}

\subsection{Sơ đồ use case tổng quát}

Hình~\ref{fig:use-case} tổng hợp các ca sử dụng của ba tác nhân đã xác định ở
mục~\ref{sec:tac-nhan}. Quan hệ kế thừa quyền giữa các tác nhân được thể hiện
bằng việc tác nhân cấp cao dùng được mọi ca của cấp thấp hơn.

% So do use case tong quat.
% Bo tri MOT cot ellipse va dat moi tac nhan ngang tam khoi ca su dung cua
% rieng no: nho vay khong duong noi nao cat qua ellipse khac.
\begin{figure}[H]
\centering
\adjustbox{max totalsize={\textwidth}{0.84\textheight}}{%
\begin{tikzpicture}[x=1cm, y=1cm,
    uc/.style={ucell, minimum width=7.6cm, minimum height=1.05cm, text width=6.0cm, align=center},
    inh/.style={thick, -{Triangle[open, length=3mm, width=3mm]}}]

  % ---------- Cac ca su dung, mot cot ----------
  \node[uc] (uc1)  at (9.5,   0.0) {UC1. Duyệt và tìm kiếm món ăn};
  \node[uc] (uc2)  at (9.5,  -1.5) {UC2. Lọc theo danh mục, vùng miền,\\độ khó, thời gian nấu};
  \node[uc] (uc3)  at (9.5,  -3.0) {UC3. Sắp xếp danh sách món};
  \node[uc] (uc4)  at (9.5,  -4.5) {UC4. Xem chi tiết món ăn};
  \node[uc] (uc5)  at (9.5,  -6.0) {UC5. Gợi ý món theo nguyên liệu sẵn có};
  \node[uc] (uc6)  at (9.5,  -7.5) {UC6. Đăng ký tài khoản};
  \node[uc] (uc7)  at (9.5,  -9.0) {UC7. Đăng nhập và đăng xuất};

  \node[uc] (uc8)  at (9.5, -11.0) {UC8. Sinh thực đơn tuần tự động};
  \node[uc] (uc9)  at (9.5, -12.5) {UC9. Chỉnh sửa thủ công từng bữa};
  \node[uc] (uc10) at (9.5, -14.0) {UC10. Tạo danh sách đi chợ};
  \node[uc] (uc11) at (9.5, -15.5) {UC11. Đánh dấu đã mua};
  \node[uc] (uc12) at (9.5, -17.0) {UC12. Quản lý món yêu thích};

  \node[uc] (uc13) at (9.5, -19.0) {UC13. Quản lý món ăn};
  \node[uc] (uc14) at (9.5, -20.5) {UC14. Quản lý nguyên liệu};
  \node[uc] (uc15) at (9.5, -22.0) {UC15. Quản lý danh mục};
  \node[uc] (uc16) at (9.5, -23.5) {UC16. Xem thống kê tổng quan};

  % ---------- Bien he thong ----------
  \node[font=\bfseries] (systitle) at (9.5, 1.5) {Hệ thống CookingAdvisor};
  \node[ucsys, fit=(uc1)(uc16)(systitle), inner sep=8pt] {};

  % ---------- Tac nhan, dat ngang tam khoi ca su dung cua minh ----------
  \stickman{khach}{1.6}{-4.5}{Khách}
  \stickman{nguoidung}{1.6}{-14.0}{Người dùng}
  \stickman{quantri}{1.6}{-21.2}{Quản trị viên}

  \foreach \t in {uc1,uc2,uc3,uc4,uc5,uc6,uc7}
    \draw[ucline] (khach) -- (\t.west);
  \foreach \t in {uc8,uc9,uc10,uc11,uc12}
    \draw[ucline] (nguoidung) -- (\t.west);
  \foreach \t in {uc13,uc14,uc15,uc16}
    \draw[ucline] (quantri) -- (\t.west);

  % ---------- Ke thua quyen giua cac tac nhan ----------
  \draw[inh] (nguoidung) -- ++(-2.4,0) |- (khach);
  \node[lbl, rotate=90, anchor=south] at (-0.8,-9.2) {kế thừa quyền};
  \draw[inh] (quantri) -- ++(-2.4,0) |- (nguoidung);
  \node[lbl, rotate=90, anchor=south] at (-0.8,-17.6) {kế thừa quyền};

\end{tikzpicture}%
}
\caption{Sơ đồ use case tổng quát của hệ thống}
\label{fig:use-case}
\end{figure}


\subsection{Đặc tả các ca sử dụng}

Bảng~\ref{tab:uc-tom-tat} tóm tắt mười một ca sử dụng của hệ thống kèm tác nhân
và điểm xử lý đáng chú ý của từng ca. Ba ca ứng với ba chức năng lõi (UC01,
UC02 và UC04) được đặc tả đầy đủ theo mẫu tác nhân, mục tiêu, tiền điều kiện,
luồng chính, luồng thay thế và hậu điều kiện ở Phụ lục E.


\begin{longtable}{|P{1.2cm}|P{3.3cm}|P{2.4cm}|P{6.2cm}|}
\caption{Tóm tắt các ca sử dụng của hệ thống}
\label{tab:uc-tom-tat}\\
\hline
\bfseries Mã & \bfseries Tên ca & \bfseries Tác nhân & \bfseries Điểm xử lý
  đáng chú ý \\ \hline
\endfirsthead
\hline
\bfseries Mã & \bfseries Tên ca & \bfseries Tác nhân & \bfseries Điểm xử lý
  đáng chú ý \\ \hline
\endhead
\endfoot
UC01 & Gợi ý món theo nguyên liệu sẵn có & Cả ba tác nhân & Tính độ phủ rồi xếp
  hạng theo bốn tiêu chí; chưa chọn nguyên liệu nào thì giữ trạng thái ban đầu
  kèm hướng dẫn \\ \hline
UC02 & Sinh thực đơn tuần tự động & Người dùng & Sinh đủ 21 suất trong một thao
  tác rồi lưu một lần; không có món nào thoả vùng miền thì ném ngoại lệ \\ \hline
UC03 & Chỉnh sửa thủ công một suất ăn & Người dùng & Đổi hoặc xóa món trong một
  ô, tổng năng lượng ngày được tính lại; ngày trống cả ba bữa vẫn giữ cột \\ \hline
UC04 & Tạo danh sách đi chợ & Người dùng & Nạp thực đơn kèm điều kiện lọc theo
  tài khoản, gộp nhóm theo cặp nguyên liệu và đơn vị rồi cộng dồn; đã có danh
  sách thì ghi đè \\ \hline
UC05 & Đánh dấu đã mua & Người dùng & Tick từng dòng nguyên liệu, hệ thống đảo
  trạng thái \code{IsPurchased}, gạch ngang dòng và cập nhật thanh tiến
  độ \\ \hline
UC06 & Duyệt, tìm kiếm, lọc và sắp xếp món ăn & Cả ba tác nhân & Mọi điều kiện
  ghép vào một truy vấn, phân trang 9 món; tham số sắp xếp không hợp lệ được đưa
  về mặc định \\ \hline
UC07 & Xem chi tiết món ăn & Cả ba tác nhân & Hiển thị thông số, nguyên liệu
  định lượng và các bước nấu; người dùng đã đăng nhập thấy thêm nút yêu
  thích \\ \hline
UC08 & Đăng ký tài khoản & Khách & Kiểm tra định dạng thư điện tử, mật khẩu tối
  thiểu 8 ký tự và khớp xác nhận; tài khoản mới được gán vai trò \code{User} và
  đăng nhập luôn \\ \hline
UC09 & Đăng nhập, đăng xuất & Cả ba tác nhân & Khi thông tin sai, hệ thống báo
  lỗi chung mà không tiết lộ sai ở trường nào \\ \hline
UC10 & Quản lý món ăn & Quản trị viên & Nhập thông số món và gán từng nguyên
  liệu kèm khối lượng, đơn vị; xóa món đang nằm trong thực đơn của người dùng bị
  chặn theo ràng buộc Restrict ở Bảng~\ref{tab:hanh-vi-xoa} \\ \hline
UC11 & Quản lý nguyên liệu và danh mục & Quản trị viên & Xóa nguyên liệu đang
  được công thức hoặc danh sách đi chợ tham chiếu, hay danh mục đang có món sử
  dụng, đều bị chặn kèm thông báo giải thích \\ \hline
\end{longtable}

\section{Kiến trúc hệ thống}
\label{sec:kien-truc}

\subsection{Kiến trúc phân tầng}

Hệ thống được tổ chức theo mô hình phân tầng, cụ thể hóa mẫu MVC đã trình bày ở
mục~\ref{sec:mvc}. Điểm khác biệt so với mẫu MVC cơ bản là việc bổ sung
\textbf{tầng Service} nằm giữa Controller và tầng truy cập dữ liệu.

% Kien truc phan tang: 5 tang xep doc; luong yeu cau di xuong ve o canh trai,
% luong tra ve di len o canh phai; khoi Identity dat ben phai, noi net dut.
\begin{figure}[H]
\centering
\resizebox{\textwidth}{!}{%
\begin{tikzpicture}[x=1cm, y=1cm]

  \node[tier, minimum width=9.6cm] (B)  at (0,0)
    {Trình duyệt\\(Razor View + Bootstrap)};
  \node[tier, minimum width=9.6cm] (C)  at (0,-3.0)
    {Tầng Controllers\\Home, Recipe, Suggestion, Menu,\\ShoppingList, Favorite, Account, Admin};
  \node[tier, minimum width=9.6cm] (S)  at (0,-6.4)
    {Tầng Services\\RecipeService, SuggestionService,\\MenuPlannerService, ShoppingListService,\\FavoriteService};
  \node[tier, minimum width=9.6cm] (D)  at (0,-9.8)
    {AppDbContext\\(Entity Framework Core)};
  \node[store, minimum width=6.0cm, minimum height=1.7cm] (DB) at (0,-12.3)
    {SQL Server};

  % --- Luong yeu cau: di xuong, ve o canh trai ---
  \draw[fl] ([xshift=-3.2cm]B.south) -- node[lbl, left] {HTTP request}
            ([xshift=-3.2cm]C.north);
  \draw[fl] ([xshift=-3.2cm]C.south) -- node[lbl, left] {gọi nghiệp vụ}
            ([xshift=-3.2cm]S.north);
  \draw[fl] ([xshift=-3.2cm]S.south) -- node[lbl, left] {LINQ}
            ([xshift=-3.2cm]D.north);
  \draw[fl] ([xshift=-1.8cm]D.south) -- node[lbl, left] {SQL}
            ([xshift=-1.8cm]DB.north);

  % --- Luong tra ve: di len, ve o canh phai ---
  \draw[fl] ([xshift=1.8cm]DB.north) -- node[lbl, right] {bản ghi}
            ([xshift=1.8cm]D.south);
  \draw[fl] ([xshift=3.2cm]D.north)  -- node[lbl, right] {thực thể}
            ([xshift=3.2cm]S.south);
  \draw[fl] ([xshift=3.2cm]S.north)  -- node[lbl, right] {ViewModel}
            ([xshift=3.2cm]C.south);
  \draw[fl] ([xshift=3.2cm]C.north)  -- node[lbl, right] {HTML}
            ([xshift=3.2cm]B.south);

  % --- ASP.NET Core Identity, cat ngang qua tang Controllers va DbContext ---
  \node[tier, minimum width=5.4cm, minimum height=1.7cm] (I) at (9.4,-6.4)
    {ASP.NET Core Identity\\xác thực và phân quyền};
  \draw[arrd] (I.north) |- node[lbl, pos=0.82, above] {bộ lọc [Authorize]}
            ([yshift=0.35cm]C.east);
  \draw[arrd] (I.south) |- node[lbl, pos=0.82, below] {bảng AspNet*}
            ([yshift=-0.35cm]D.east);

\end{tikzpicture}%
}
\caption{Kiến trúc phân tầng và luồng xử lý một yêu cầu}
\label{fig:kien-truc}
\end{figure}


Lý do tách tầng Service là để logic nghiệp vụ không phụ thuộc vào hạ tầng web.
Nhờ đó ba thuật toán lõi có thể được kiểm thử bằng kiểm thử đơn vị mà không cần
khởi động máy chủ web hay mô phỏng đối tượng \code{HttpContext}. Kết quả cụ thể
của lựa chọn này được trình bày ở Chương~\ref{ch:ket-qua}.

\subsection{Tổ chức mã nguồn}

Mã nguồn được chia theo trách nhiệm: \code{Controllers} tiếp nhận yêu cầu và gọi
Service, \code{Areas/Admin} tách riêng khu quản trị, \code{Services} chứa ba
thuật toán lõi, \code{Models} và \code{Data} giữ thực thể cùng
\code{AppDbContext}, \code{ViewModels} là dữ liệu truyền cho \code{Views}. Cây
thư mục đầy đủ nằm ở Phụ lục B.

Nguyên tắc phân chia là mỗi thư mục ứng với đúng một trách nhiệm, và chiều phụ
thuộc luôn đi từ ngoài vào trong: View phụ thuộc ViewModel, Controller phụ thuộc
Service, Service phụ thuộc \code{AppDbContext} và Models. Không có chiều ngược
lại, nhờ đó có thể thay giao diện mà không đụng tới thuật toán.

\subsection{Mô hình triển khai}

Hình~\ref{fig:trien-khai} mô tả cách hệ thống được triển khai. Ứng dụng web chạy
trên máy chủ Kestrel của .NET và kết nối tới một máy chủ SQL Server riêng biệt.
Như đã phân tích ở mục~\ref{sec:sqlserver}, máy chủ cơ sở dữ liệu có thể chạy
trong container hoặc cài trực tiếp; mã nguồn không phân biệt hai trường hợp này.

% So do trien khai: 3 nut thiet bi xep doc, moi nut chua cac thanh phan
% duoc trien khai; khung net dut ghi chu hai phuong an chay SQL Server.
\begin{figure}[H]
\centering
\resizebox{\textwidth}{!}{%
\begin{tikzpicture}[x=1cm, y=1cm,
    % Khung nut thiet bi ve SAU cac thanh phan ben trong nen phai fill=none,
    % neu khong khung se de len va che mat chu.
    devnode/.style={draw, thick, fill=none, inner sep=5pt}]

  % ---------- Nut 1: may nguoi dung ----------
  \node[artifact, minimum width=5.6cm] (a1) at (0,0) {Trình duyệt web (Chrome)};
  \node[above=2mm of a1, font=\small\bfseries] (t1) {Máy người dùng};
  \node[devnode, fit=(a1)(t1)] (n1) {};

  % ---------- Nut 2: may chu ung dung ----------
  \node[artifact, minimum width=6.6cm] (b1) at (0,-3.6)
    {CookingAdvisor.dll (ASP.NET Core MVC)};
  \node[artifact, minimum width=6.6cm, below=1.5mm of b1] (b2)
    {Controllers, Services, ViewModels};
  \node[artifact, minimum width=6.6cm, below=1.5mm of b2] (b3)
    {Views (Razor)};
  \node[artifact, minimum width=6.6cm, below=1.5mm of b3] (b4)
    {wwwroot: CSS, JavaScript, ảnh};
  \node[above=2mm of b1, font=\small\bfseries, align=center] (t2)
    {Máy chủ ứng dụng\\{\scriptsize\mdseries .NET 10, máy chủ Kestrel}};
  \node[devnode, fit=(b1)(b2)(b3)(b4)(t2)] (n2) {};

  % ---------- Nut 3: may chu co so du lieu ----------
  \node[artifact, minimum width=5.6cm] (c1) at (0,-9.4) {SQL Server 2022};
  \node[artifact, minimum width=5.6cm, below=1.5mm of c1] (c2)
    {Cơ sở dữ liệu\\CookingAdvisor};
  \node[above=2mm of c1, font=\small\bfseries] (t3) {Máy chủ cơ sở dữ liệu};
  \node[devnode, fit=(c1)(c2)(t3)] (n3) {};

  % ---------- Ghi chu hai phuong an ----------
  \node[devgroup, right=1.4cm of n3, text width=6.4cm, align=left,
        font=\scriptsize] (note)
    {Hai phương án tương đương:\\[2pt]
     1. Container Docker hoặc OrbStack trên macOS.\\[2pt]
     2. SQL Server Express hoặc Developer cài trực tiếp trên Windows.};
  \draw[thick, dashed] (n3.east) -- (note.west);

  % ---------- Ket noi ----------
  \draw[arr] (n1.south) -- node[lbl, right] {HTTP} (n2.north);
  \draw[arr] (n2.south) -- node[lbl, right, align=left]
    {TDS, cổng 1433\\{\scriptsize chuỗi kết nối lấy từ User Secrets}}
    (n3.north);

\end{tikzpicture}%
}
\caption{Sơ đồ triển khai hệ thống}
\label{fig:trien-khai}
\end{figure}


\section{Thiết kế lớp}
\label{sec:thiet-ke-lop}

Tầng thực thể ánh xạ một - một với các bảng đã mô tả ở mục~\ref{sec:mo-ta-bang};
ngoài các thuộc tính vô hướng, mỗi lớp thực thể còn có các thuộc tính điều hướng
để EF Core dựng quan hệ, nên không cần sơ đồ riêng ngoài sơ đồ quan hệ thực thể
ở Hình~\ref{fig:erd}. Hình~\ref{fig:lop-service} thể hiện sơ đồ lớp của tầng
dịch vụ và quan hệ với tầng điều khiển.

\begin{figure}[H]
\centering
\resizebox{\textwidth}{!}{%
\begin{tikzpicture}[node distance=1.4cm and 1.4cm]

  % --- Hang 1: Controller --------------------------------------
  \node[procs] (RC) at (0,11) {RecipeController};
  \node[procs] (SC) at (8,11) {SuggestionController};
  \node[procs] (MC) at (16,11) {MenuController};
  \node[procs] (SLC) at (24,11) {ShoppingListController};
  \node[procs] (FC) at (32,11) {FavoriteController};

  % --- Hang 2: Service -------------------------------------------
  \node[classbox] (RS) at (0,5) {\textbf{RecipeService}
    \nodepart{two} +PageSize: int
    \nodepart{three} +SearchRecipesAsync(RecipeFilterViewModel)\\+GetHomeAsync(int)\\+GetRecipeDetailAsync(int)};
  \node[classbox] (SS) at (8,5) {\textbf{SuggestionService}
    \nodepart{two} ~
    \nodepart{three} +GetIngredientOptionsAsync()\\+SuggestAsync(IReadOnlyCollection)};
  \node[classbox] (MPS) at (16,5) {\textbf{MenuPlannerService}
    \nodepart{two} -DailyCalorieTarget: int\\-Meals: array
    \nodepart{three} +GenerateWeeklyPlanAsync(userId,\\planName, weekStartDate, region)};
  \node[classbox] (SLS) at (24,5) {\textbf{ShoppingListService}
    \nodepart{two} ~
    \nodepart{three} +GenerateFromMenuPlanAsync(userId, menuPlanId)};
  \node[classbox] (FS) at (32,5) {\textbf{FavoriteService}
    \nodepart{two} ~
    \nodepart{three} +GetFavoritesAsync(userId)\\+IsFavoriteAsync(userId, recipeId)\\+ToggleFavoriteAsync(userId, recipeId)};

  % --- Hang 3: AppDbContext ---------------------------------------
  \node[classbox] (DB) at (16,-2.5) {\textbf{AppDbContext}
    \nodepart{two} +DbSet$<$Recipe$>$ Recipes\\+DbSet$<$Ingredient$>$ Ingredients\\+DbSet$<$Category$>$ Categories\\+DbSet$<$RecipeIngredient$>$ RecipeIngredients\\+DbSet$<$MenuPlan$>$ MenuPlans\\+DbSet$<$MenuPlanItem$>$ MenuPlanItems\\+DbSet$<$Favorite$>$ Favorites\\+DbSet$<$ShoppingList$>$ ShoppingLists\\+DbSet$<$ShoppingListItem$>$ ShoppingListItems
    \nodepart{three} \#OnModelCreating(ModelBuilder)};

  % --- Quan he: Controller -> Service (net dut) -------------------
  \draw[arrd] (RC) -- (RS) node[lbl, pos=0.5] {sử dụng};
  \draw[arrd] (RC.south) -- ++(0,-2.7) -| (FS.west) node[lbl, pos=0.15] {sử dụng};
  \draw[arrd] (SC) -- (SS) node[lbl, pos=0.5] {sử dụng};
  \draw[arrd] (MC) -- (MPS) node[lbl, pos=0.5] {sử dụng};
  \draw[arrd] (SLC) -- (SLS) node[lbl, pos=0.5] {sử dụng};
  \draw[arrd] (FC) -- (FS) node[lbl, pos=0.5] {sử dụng};

  % --- Quan he: Service -> AppDbContext (net dut) ------------------
  \draw[arrd] (RS) -- (DB);
  \draw[arrd] (SS) -- (DB);
  \draw[arrd] (MPS) -- (DB);
  \draw[arrd] (SLS) -- (DB);
  \draw[arrd] (FS) -- (DB);

\end{tikzpicture}%
}
\caption{Sơ đồ lớp tầng dịch vụ và quan hệ với tầng điều khiển}
\label{fig:lop-service}
\end{figure}


Bảng~\ref{tab:trach-nhiem-service} tóm tắt trách nhiệm của từng lớp dịch vụ.

\begin{table}[H]
\centering
\caption{Trách nhiệm của các lớp dịch vụ}
\label{tab:trach-nhiem-service}
\begin{tabularx}{\textwidth}{|P{5.7cm}|Y|}
\hline
\bfseries Lớp & \bfseries Trách nhiệm chính \\ \hline
\code{RecipeService} & Tìm kiếm, lọc, sắp xếp, phân trang; lấy chi tiết món; dữ
  liệu trang chủ \\ \hline
\code{SuggestionService} & Thuật toán gợi ý theo nguyên liệu \\ \hline
\code{MenuPlannerService} & Thuật toán sinh thực đơn tuần \\ \hline
\code{ShoppingListService} & Thuật toán gộp nhóm sinh danh sách đi chợ \\ \hline
\code{FavoriteService} & Thêm, bỏ, kiểm tra và liệt kê món yêu thích \\ \hline
\end{tabularx}
\end{table}

Các lớp dịch vụ được đăng ký vào bộ chứa tiêm phụ thuộc với vòng đời
\code{Scoped}, tương ứng với vòng đời của \code{AppDbContext}. Việc thống nhất
vòng đời là bắt buộc: nếu đăng ký Service ở vòng đời dài hơn
\code{AppDbContext}, Service sẽ giữ lại một ngữ cảnh đã bị hủy và phát sinh lỗi
khi có yêu cầu tiếp theo.

\section{Thiết kế luồng xử lý}
\label{sec:luong-xu-ly}

Mỗi chức năng lõi được mô tả bằng một sơ đồ tuần tự, cho thấy thứ tự trao đổi
thông điệp giữa các thành phần bên trong hệ thống; các nhánh lỗi và bước ra
quyết định bên trong từng thuật toán được thể hiện đầy đủ ở các lưu đồ trong
mục~\ref{sec:cai-dat-thuat-toan}.

\subsection{Luồng gợi ý theo nguyên liệu}

\begin{figure}[H]
\centering
\adjustbox{max totalsize={\textwidth}{0.82\textheight}}{%
\begin{tikzpicture}[x=3.4cm, y=1.15cm]

  % --- Doi tuong tren dau ------------------------------------------
  \node[seqpart] (ND) at (0,0) {Người dùng};
  \node[seqpart] (VW) at (1,0) {View\\Suggestion/Index};
  \node[seqpart] (CT) at (2,0) {Suggestion\\Controller};
  \node[seqpart] (SV) at (3,0) {Suggestion\\Service};
  \node[seqpart] (DB) at (4,0) {AppDbContext /\\SQL Server};

  \foreach \x in {0,1,2,3,4} \draw[seqlife] (\x,-0.5) -- (\x,-21.6);

  % --- Cac thong diep chinh -----------------------------------------
  \draw[seqmsg] (0,-1.4) -- node[above=2pt, fill=white, inner sep=2pt, font=\scriptsize]{1. Tick nguyên liệu đang có} (1,-1.4);
  \draw[seqmsg] (0,-2.3) -- node[above=2pt, fill=white, inner sep=2pt, font=\scriptsize]{2. Bấm ``Gợi ý món''} (1,-2.3);
  \draw[seqmsg] (1,-3.2) -- node[above=2pt, fill=white, inner sep=2pt, font=\scriptsize]{3. GET /Suggestion?ingredientIds=...} (2,-3.2);
  \draw[seqmsg] (2,-4.1) -- node[above=2pt, fill=white, inner sep=2pt, font=\scriptsize]{4. GetIngredientOptionsAsync()} (3,-4.1);
  \draw[seqmsg] (3,-5.0) -- node[above=2pt, fill=white, inner sep=2pt, font=\scriptsize]{5. Truy vấn Ingredients, sắp theo tên} (4,-5.0);
  \draw[seqret] (4,-5.9) -- node[above=2pt, fill=white, inner sep=2pt, font=\scriptsize]{6. Danh sách nguyên liệu} (3,-5.9);
  \draw[seqret] (3,-6.8) -- node[above=2pt, fill=white, inner sep=2pt, font=\scriptsize]{7. List IngredientOptionViewModel} (2,-6.8);
  \draw[seqmsg] (2,-7.7) -- node[above=2pt, fill=white, inner sep=2pt, font=\scriptsize]{8. SuggestAsync(ownedIngredientIds)} (3,-7.7);

  % --- Khung alt ------------------------------------------------------
  \draw[seqframe] (1.5,-8.2) rectangle (4.55,-18.45);
  \node[seqlab, anchor=north west] at (1.5,-8.2) {alt};
  \node[font=\scriptsize\itshape, anchor=north west] at (1.65,-8.55) {Tập nguyên liệu rỗng};

  \draw[seqret] (3,-9.5) -- node[above=2pt, fill=white, inner sep=2pt, font=\scriptsize]{9. Danh sách rỗng} (2,-9.5);

  \draw[seqframe, dashed] (1.5,-10.0) -- (4.55,-10.0);
  \node[font=\scriptsize\itshape, anchor=north west] at (1.65,-10.35) {Có ít nhất một nguyên liệu};

  \draw[seqmsg] (3,-11.3) -- node[above=2pt, fill=white, inner sep=2pt, font=\scriptsize]{10. Lấy món có giao với tập đã chọn} (4,-11.3);
  \draw[seqret] (4,-12.2) -- node[above=2pt, fill=white, inner sep=2pt, font=\scriptsize]{11. Món kèm danh sách nguyên liệu} (3,-12.2);

  % --- Khung loop long ben trong alt ---------------------------------
  \draw[seqframe] (2.55,-12.7) rectangle (3.85,-16.35);
  \node[seqlab, anchor=north west] at (2.55,-12.7) {loop};
  \node[font=\scriptsize\itshape, anchor=north west] at (2.85,-13.0) {Với mỗi món R};

  \draw[seqmsg] (3,-13.6) -- (3.35,-13.6) -- (3.35,-13.95) -- (3,-13.95);
  \node[fill=white, inner sep=2pt, font=\scriptsize, align=center, anchor=west] at (3.4,-13.775) {12. missing =\\$I_R \setminus A$};

  \draw[seqmsg] (3,-14.6) -- (3.35,-14.6) -- (3.35,-14.95) -- (3,-14.95);
  \node[fill=white, inner sep=2pt, font=\scriptsize, align=center, anchor=west] at (3.4,-14.775) {13. matched =\\$|I_R| - |\text{missing}|$};

  \draw[seqmsg] (3,-15.6) -- (3.35,-15.6) -- (3.35,-15.95) -- (3,-15.95);
  \node[fill=white, inner sep=2pt, font=\scriptsize, align=center, anchor=west] at (3.4,-15.775) {14. coverage =\\matched / $|I_R|$};

  % --- Sau loop, van trong nhanh 2 cua alt ---------------------------
  \draw[seqmsg] (3,-16.85) -- (3.35,-16.85) -- (3.35,-17.2) -- (3,-17.2);
  \node[fill=white, inner sep=2pt, font=\scriptsize, align=center, anchor=west] at (3.4,-17.025) {15. Xếp hạng\\bốn tiêu chí};

  \draw[seqret] (3,-18.0) -- node[above=2pt, fill=white, inner sep=2pt, font=\scriptsize]{16. Danh sách đã xếp hạng} (2,-18.0);

  % --- Ket thuc alt, tra ve cho nguoi dung ----------------------------
  \draw[seqret] (2,-19.25) -- node[above=2pt, fill=white, inner sep=2pt, font=\scriptsize]{17. SuggestionViewModel} (1,-19.25);
  \draw[seqret] (1,-21.1) -- node[above=2pt, fill=white, inner sep=2pt, font=\scriptsize, text width=4.6cm, align=center]{18. Lưới món kèm nhãn ``Đủ nguyên liệu'' hoặc ``Thiếu N nguyên liệu''} (0,-21.1);

\end{tikzpicture}%
}
\caption{Sơ đồ tuần tự chức năng gợi ý theo nguyên liệu}
\label{fig:seq-goi-y}
\end{figure}


Điểm cần chú ý ở luồng này là hai lời gọi tách biệt tới tầng dịch vụ: một lời gọi
lấy danh sách nguyên liệu để dựng biểu mẫu và một lời gọi thực hiện thuật toán
gợi ý. Nhờ tách như vậy, khi người dùng chưa chọn nguyên liệu nào, hệ thống vẫn
dựng được biểu mẫu đầy đủ mà không phải chạy thuật toán.

\subsection{Luồng sinh thực đơn tuần}

Luồng sinh thực đơn đi theo cùng trình tự: Controller nhận biểu mẫu, gọi
\code{MenuPlannerService}, tầng dịch vụ nạp danh sách món ứng viên rồi chạy vòng
lặp hai mức tương ứng với 7 ngày và 3 bữa mỗi ngày. Toàn bộ 21 suất được sinh
trong bộ nhớ rồi mới lưu xuống cơ sở dữ liệu một lần, thay vì lưu từng suất, nhờ
đó tránh được 21 lần đi vòng tới máy chủ cơ sở dữ liệu.

\subsection{Luồng sinh danh sách đi chợ}

Luồng này đi theo đúng trình tự của hai luồng trên: Controller nhận yêu cầu, gọi
\code{ShoppingListService}, tầng dịch vụ nạp thực đơn kèm toàn bộ nguyên liệu
rồi ghi danh sách đã gộp xuống cơ sở dữ liệu. Điểm đáng chú ý ở luồng này là bước nạp dữ liệu có kèm điều kiện lọc theo
\code{UserId}. Nhờ đó, nếu người dùng cố tình gửi mã thực đơn của người khác,
truy vấn trả về rỗng và hệ thống ném ngoại lệ, thay vì sinh danh sách đi chợ từ
dữ liệu không thuộc quyền của họ. Đây là cách hiện thực yêu cầu phi chức năng N3
đã nêu ở mục~\ref{sec:yeu-cau-phi-chuc-nang}.

\section{Cài đặt các thuật toán lõi}
\label{sec:cai-dat-thuat-toan}

Mã giả đầy đủ của cả ba thuật toán, có đánh số dòng, được trình bày ở Phụ lục H.
Phần dưới đây giải thích các quyết định cài đặt quan trọng của từng thuật toán,
kèm số dòng tương ứng trong mã giả đó.

\subsection{Thuật toán gợi ý theo nguyên liệu}
\label{sec:thuat-toan-goi-y}

\textit{Lọc sớm và tra cứu bằng bảng băm (dòng 5-6).} Điều kiện ``có ít nhất một
nguyên liệu thuộc A'' được đẩy xuống câu truy vấn thay vì lọc trong bộ nhớ:
những món không chung nguyên liệu nào chắc chắn có $\mathrm{coverage} = 0$,
không đáng được gợi ý, nên loại bỏ sớm giảm khối lượng dữ liệu chuyển lên tầng
ứng dụng. Tập $A$ được chuyển thành bảng băm vì phép kiểm tra thuộc tập ở dòng
11 lặp lại $|I_r|$ lần cho mỗi món; chi phí mỗi phép kiểm tra giảm từ $O(|A|)$
xuống trung bình $O(1)$.

\textit{Bảo đảm thứ tự tất định (dòng 12 và 23).} Danh sách nguyên liệu thiếu
được sắp theo bảng chữ cái để nhãn ``Thiếu: X, Y'' hiển thị ổn định giữa các lần
chạy; khóa sắp xếp cuối cùng theo tên món, như đã lập luận ở
mục~\ref{sec:do-do}, bảo đảm kết quả có thứ tự tất định, thuận tiện cho kiểm thử
tự động.

\noindent\textbf{Độ phức tạp:} $O(nk + n\log n)$ với $n$ là số món thỏa điều kiện
lọc và $k$ là số nguyên liệu trung bình mỗi món.

\subsection{Thuật toán sinh thực đơn tuần}
\label{sec:thuat-toan-thuc-don}

\textit{Dòng 1 và 19-20, thứ tự nới lỏng ràng buộc.} Ràng buộc vùng miền được áp
ngay từ đầu và không bao giờ được nới. Hai ràng buộc còn lại được nới theo thứ
tự: ``không lặp'' trước, ``hợp bữa'' sau. Lý do là việc lặp lại một món trong
tuần gây khó chịu ít hơn so với việc đề xuất một món hoàn toàn không hợp bữa,
chẳng hạn đề xuất món tráng miệng cho bữa sáng.

\textit{Dòng 13, mốc năng lượng lũy tiến theo bữa.} Giá trị \code{mealTarget}
được tính lũy tiến theo chỉ số bữa, bằng $\frac{2000 \times (m+1)}{3}$, cho các
mốc lần lượt là 667, 1333 và 2000 kcal. Cách tính này khiến tiêu chí ở dòng 30 so
sánh \textbf{tổng năng lượng tích lũy tới hết bữa hiện tại} với mốc mong đợi,
thay vì so sánh riêng từng bữa một cách độc lập.

\textit{Dòng 22-27, hai chế độ sắp xếp khác nhau.} Đây là điểm cài đặt quan
trọng nhất và cũng là nơi phát sinh khiếm khuyết đã nêu ở
mục~\ref{sec:tham-lam}. Khi còn món chưa dùng (dòng 26), tiêu chí đầu tiên là ưu
tiên món yêu thích; nhưng khi buộc phải lặp (dòng 23-24), tiêu chí đầu tiên phải
chuyển thành \textbf{số lần đã dùng ít nhất}. Nếu bỏ qua sự phân biệt này, trong
nhánh lặp tiêu chí quyết định thực tế lùi về khoảng cách năng lượng ở dòng 30;
do trạng thái \code{dayCalories} lặp lại theo chu kỳ giữa các ngày, cùng một món
luôn cho khoảng cách nhỏ nhất và được chọn lại liên tục. Kết quả đo được trên bộ
dữ liệu thử là một món chiếm 12 trong số các suất còn trống trong khi nhiều món
khác không được dùng lần nào.

\textit{Dòng 31, khóa sắp xếp cuối.} Sắp theo \code{Id} tăng dần để bảo đảm với
cùng đầu vào, thuật toán luôn sinh ra cùng một thực đơn; tính chất tất định này
là điều kiện cần để viết kiểm thử tự động.

\noindent\textbf{Độ phức tạp:} $O(S \cdot n \log n)$ với $S = 21$ suất và $n$ số
món ứng viên.

\subsection{Thuật toán sinh danh sách đi chợ}
\label{sec:thuat-toan-di-cho}

\textit{Dòng 2, lọc theo \code{UserId} ngay trong truy vấn.} Việc kiểm tra quyền
sở hữu được thực hiện ở tầng truy vấn chứ không phải bằng câu lệnh điều kiện sau
khi đã nạp dữ liệu; nhờ đó không tồn tại nhánh nào có thể vô tình bỏ qua bước
kiểm tra.

\textit{Dòng 6 và 10, duyệt theo suất ăn với khóa gộp nhóm ghép cặp.} Nếu duyệt
theo tập món khác nhau, một món xuất hiện ba lần trong tuần chỉ được tính nguyên
liệu một lần, dẫn tới mua thiếu. Đơn vị đo phải nằm trong khóa gộp nhóm để không
cộng dồn hai giá trị khác đơn vị, như đã phân tích ở mục~\ref{sec:gop-nhom}.

\textit{Dòng 16-21, ghi đè thay vì tạo mới.} Do ràng buộc duy nhất trên
\code{MenuPlanId} (mục~\ref{sec:mo-ta-bang}), mỗi thực đơn chỉ có tối đa một danh
sách đi chợ; khi người dùng sửa thực đơn rồi tạo lại danh sách, hệ thống xóa các
dòng cũ và ghi lại từ đầu.

\noindent\textbf{Độ phức tạp:} $O(S \cdot k)$ với $S$ là số suất ăn và $k$ là số
nguyên liệu trung bình mỗi món.

\section{Một số điểm cài đặt khác}
\label{sec:cai-dat-khac}

\subsection{Tìm kiếm, lọc và sắp xếp}

Toàn bộ tham số tìm kiếm được gom vào một lớp \code{RecipeFilterViewModel} và ánh
xạ tự động nhờ cơ chế model binding. Các điều kiện lọc được ghép dần vào đối
tượng \code{IQueryable}, do đó chỉ một câu lệnh SQL duy nhất được sinh ra dù
người dùng chọn bao nhiêu điều kiện.

Tùy chọn sắp xếp được biểu diễn bằng kiểu liệt kê \code{RecipeSortOrder}. Việc
dùng kiểu liệt kê thay vì chuỗi ký tự có ý nghĩa về an toàn: \textbf{bản thân
kiểu liệt kê đóng vai trò danh sách trắng}, giá trị không hợp lệ trên thanh địa
chỉ sẽ được ánh xạ về giá trị mặc định, nên không có chuỗi nào do người dùng nhập
vào có thể đi tới biểu thức sắp xếp.

Mọi phương án sắp xếp đều kết thúc bằng khóa phụ là tên món. Thiếu khóa này, các
món có cùng giá trị sắp xếp sẽ không có thứ tự tất định và một món có thể xuất
hiện ở hai trang khác nhau khi phân trang. Đây chính là khiếm khuyết thứ hai được
trình bày ở mục~\ref{sec:khiem-khuyet}.

\subsection{Bảo mật}

Sáu nguy cơ thường gặp đều được chặn ngay trong thiết kế. Nguy cơ \textbf{chèn mã
SQL} không tồn tại vì mọi truy vấn đi qua LINQ của EF Core và được tham số hóa tự
động. Nguy cơ \textbf{lộ mật khẩu} được giảm thiểu vì Identity chỉ lưu giá trị
băm có muối. Nguy cơ \textbf{giả mạo yêu cầu liên trang} được chặn bằng token
chống giả mạo trên mọi biểu mẫu POST. \textbf{Truy cập vượt quyền chức năng} bị
chặn bởi \code{[Authorize(Roles = "Admin")]} trên toàn bộ khu quản trị, còn
\textbf{truy cập vượt quyền dữ liệu} bị chặn bằng điều kiện lọc theo
\code{UserId} đặt ngay trong truy vấn. Cuối cùng, \textbf{chuỗi kết nối} nằm
trong User Secrets chứ không nằm trong mã nguồn.
